\documentclass[11pt,a4paper]{article}

% ── Packages ──────────────────────────────────────────────────────────────────
\usepackage[utf8]{inputenc}
\usepackage[T1]{fontenc}
\usepackage{lmodern}
\usepackage[margin=2.5cm]{geometry}
\usepackage{amsmath,amssymb,amsfonts,mathtools}
\usepackage{booktabs}
\usepackage{tabularx}
\usepackage{array}
\usepackage{enumitem}
\usepackage{hyperref}
\usepackage{xcolor}
\usepackage{graphicx}
\usepackage{float}
\usepackage{caption}
\usepackage{subcaption}
\usepackage{fancyvrb}
\usepackage[numbers,sort&compress]{natbib}
\usepackage{microtype}
\usepackage{titlesec}
\usepackage{tcolorbox}

% ── Formatting ────────────────────────────────────────────────────────────────
\hypersetup{
  colorlinks=true,
  linkcolor=blue!60!black,
  citecolor=green!50!black,
  urlcolor=blue!70!black,
}
\setlength{\parskip}{0.4em}
\setlength{\parindent}{0pt}

% ── Shape annotation box ────────────────────────────────────────────────────
\newtcolorbox{shapes}{
  colback=gray!8, colframe=gray!50, boxrule=0.3pt,
  left=4pt, right=4pt, top=2pt, bottom=2pt,
  fontupper=\small, before skip=2pt, after skip=6pt,
}

% ── Custom commands ───────────────────────────────────────────────────────────
\newcommand{\jvp}{\mathrm{JVP}}
\newcommand{\EMA}{\mathrm{EMA}}
\newcommand{\FID}{\mathrm{FID}}
\newcommand{\SWD}{\mathrm{SWD}}
\newcommand{\ie}{\textit{i.e.}}
\newcommand{\eg}{\textit{e.g.}}
\newcommand{\etal}{\textit{et~al.}}
\newcommand{\code}[1]{\texttt{#1}}
\DeclareMathOperator{\sg}{sg}        % stop-gradient
\DeclareMathOperator{\sigmoid}{\sigma}
\DeclareMathOperator{\diag}{diag}

% ══════════════════════════════════════════════════════════════════════════════
\begin{document}

% ── Title ─────────────────────────────────────────────────────────────────────
\begin{center}
  {\LARGE\bfseries NeuroMF: Latent MeanFlow for\\[4pt] One-Step 3D Brain MRI Synthesis}\\[18pt]
  {\large Technical Report --- Detailed Scientific Summary}\\[12pt]
  {\large Mario Pascual-Gonz\'{a}lez}\\[6pt]
  {\normalsize February 2026}\\[6pt]
  {\small\itshape Phases 0--5 complete (code); Phase~4 best model trained and evaluated;
  Phase~5 generation pipeline ready}
\end{center}

\vspace{1em}
\hrule
\vspace{0.5em}

\tableofcontents

\vspace{1em}
\hrule
\newpage

%═══════════════════════════════════════════════════════════════════════════════
\section{Motivation and Problem Statement}
\label{sec:motivation}

\subsection{Clinical Need}

Generative models for 3D brain MRI synthesis serve three clinical purposes:
(i)~data augmentation for rare pathologies with scarce training data,
(ii)~synthetic cohorts for privacy-preserving research, and
(iii)~counterfactual generation for explainability.
The critical bottleneck in existing methods is \textbf{sampling cost}:
state-of-the-art approaches (DDPM, flow matching, rectified flow)
require 5--1000 network evaluations per volume, making large-scale synthesis
prohibitively slow.

\subsection{Gap in the Literature}

No prior work has applied MeanFlow---or any 1-step flow-based model---to 3D
medical image synthesis. The closest works are summarised in
Table~\ref{tab:literature}.

\begin{table}[ht]
\centering
\caption{Comparison with related generative methods for brain MRI.}
\label{tab:literature}
\begin{tabular}{lccll}
\toprule
\textbf{Method} & \textbf{Space} & \textbf{NFE} & \textbf{Paradigm} & \textbf{Domain} \\
\midrule
MAISI-v2 \cite{maisi_v2}    & Latent & 5--50   & Rectified Flow       & 3D CT/MRI \\
MOTFM \cite{motfm}          & Pixel  & 10--50  & OT Flow Matching     & 3D Brain MRI \\
Med-DDPM \cite{medddpm}     & Pixel  & 1000    & DDPM                 & 3D Brain MRI \\
pMF \cite{pmf}              & Pixel  & 1       & Progressive MeanFlow & 2D Natural Images \\
\textbf{NeuroMF (ours)}     & \textbf{Latent} & \textbf{1} & \textbf{MeanFlow (iMF dual-head)} & \textbf{3D Brain MRI} \\
\bottomrule
\end{tabular}
\end{table}

Our work fills the intersection: \textbf{1-step $+$ latent $+$ 3D $+$ medical}.

\subsection{Approach Overview}

NeuroMF trains a MeanFlow model in the latent space of a frozen MAISI 3D VAE:
\begin{center}
\begin{BVerbatim}
Input MRI (1x192^3) --> Frozen MAISI VAE Encoder --> Latent (4x48^3)
                                                         |
                                                  Train MeanFlow
                                                         |
                                                         v
Synthetic MRI (1x192^3) <-- Frozen MAISI VAE Decoder <--'
\end{BVerbatim}
\end{center}

At inference, a \textbf{single forward pass} through the MeanFlow network
generates a complete 3D brain MRI volume.  The network learns the
\emph{average velocity} of a probability flow ODE, which by construction
encodes the entire transport from noise to data in one evaluation.


%═══════════════════════════════════════════════════════════════════════════════
\section{Theoretical Foundations}
\label{sec:theory}

%───────────────────────────────────────────────────────────────────────────────
\subsection{Notation and Tensor Dimensions}
\label{sec:notation}

Table~\ref{tab:notation} defines all symbols used in this report and their
tensor shapes as instantiated in code (\code{src/neuromf/}).  Equations are
presented \emph{per-sample} (batch index suppressed) unless stated otherwise;
grey \emph{Shapes} boxes show the batched implementation dimensions.

\begin{table}[ht]
\centering
\caption{Notation, symbols, and tensor dimensions.  All latent-space tensors
live in $\mathbb{R}^{C \times D \times H \times W}$ with $C{=}4$,
$D{=}H{=}W{=}48$.  The total latent dimensionality is
$d = C \cdot D \cdot H \cdot W = 4 \times 48^3 = 442{,}368$.}
\label{tab:notation}
\small
\begin{tabular}{llll}
\toprule
\textbf{Symbol} & \textbf{Meaning} & \textbf{Per-sample shape} & \textbf{Batched shape} \\
\midrule
$z_0$ & Clean latent (data) & $\mathbb{R}^{C \times D \times H \times W}$ & $(B, 4, 48, 48, 48)$ \\
$\varepsilon$ & Gaussian noise & $\mathbb{R}^{C \times D \times H \times W}$ & $(B, 4, 48, 48, 48)$ \\
$z_t$ & Noisy interpolation at time $t$ & $\mathbb{R}^{C \times D \times H \times W}$ & $(B, 4, 48, 48, 48)$ \\
$v_c$ & Conditional velocity (target) & $\mathbb{R}^{C \times D \times H \times W}$ & $(B, 4, 48, 48, 48)$ \\
$u_\theta$ & Average velocity (model output) & $\mathbb{R}^{C \times D \times H \times W}$ & $(B, 4, 48, 48, 48)$ \\
$\hat{x}$ & Denoised data estimate (x-pred) & $\mathbb{R}^{C \times D \times H \times W}$ & $(B, 4, 48, 48, 48)$ \\
$\hat{v}$ & v-head output (tangent predictor) & $\mathbb{R}^{C \times D \times H \times W}$ & $(B, 4, 48, 48, 48)$ \\
$\tilde{v}$ & JVP tangent vector & $\mathbb{R}^{C \times D \times H \times W}$ & $(B, 4, 48, 48, 48)$ \\
$V_\theta$ & Compound velocity & $\mathbb{R}^{C \times D \times H \times W}$ & $(B, 4, 48, 48, 48)$ \\
$\frac{du}{dt}$ & JVP output (directional deriv.) & $\mathbb{R}^{C \times D \times H \times W}$ & $(B, 4, 48, 48, 48)$ \\
\midrule
$t$ & Upper time (current) & $[0,1]$ & $(B,)$ \\
$r$ & Lower time (interval start) & $[0,1],\; r \le t$ & $(B,)$ \\
$h$ & Interval width $h \coloneqq t - r$ & $[0,1]$ & $(B,)$ \\
\midrule
$\mathbf{e}_t$ & Time embedding for $t$ & $\mathbb{R}^{d_\mathrm{emb}}$ & $(B, 256)$ \\
$\mathbf{e}_h$ & Time embedding for $h$ & $\mathbb{R}^{d_\mathrm{emb}}$ & $(B, 256)$ \\
$\mathbf{e}$ & Combined embedding $\mathbf{e}_t + \mathbf{e}_h$ & $\mathbb{R}^{d_\mathrm{emb}}$ & $(B, 256)$ \\
\midrule
$\ell_u, \ell_v$ & Per-sample raw losses & $\mathbb{R}_{\ge 0}$ & $(B,)$ \\
$w_u, w_v$ & Per-sample adaptive weights & $\mathbb{R}_{> 0}$ & $(B,)$ \\
$\mathcal{L}$ & Total scalar loss & $\mathbb{R}$ & scalar \\
\bottomrule
\end{tabular}
\end{table}

\textbf{Broadcasting convention.}  When a scalar-per-sample quantity $t \in
\mathbb{R}^B$ multiplies a 5-D tensor $z \in \mathbb{R}^{B \times C \times D
\times H \times W}$, $t$ is reshaped to $(B, 1, 1, 1, 1)$ via
\code{t.view(-1, 1, 1, 1, 1)} before element-wise multiplication.  We write
this broadcast as $t \odot z$ and denote the broadcast shape as
$\bar{t} \in \mathbb{R}^{B \times 1 \times 1 \times 1 \times 1}$.

%───────────────────────────────────────────────────────────────────────────────
\subsection{Operator Definitions}
\label{sec:operators}

We define two operators used throughout this report.

\textbf{Definition 1 (Stop-gradient).}  The stop-gradient operator
$\sg[\cdot]$ returns the value of its argument but blocks gradient
propagation during backpropagation:
\begin{equation}
  \sg[x] \coloneqq x, \qquad
  \frac{\partial}{\partial \theta}\,\sg[x(\theta)] = \mathbf{0}
  \label{eq:sg_def}
\end{equation}
In code, $\sg[\cdot]$ is implemented as \code{.detach()}, which detaches the
tensor from the computation graph.  For any differentiable function
$f(\theta, \sg[g(\theta)])$, the gradient $\nabla_\theta f$ treats $\sg[g]$
as a constant, so only the explicit dependence of $f$ on $\theta$ is
differentiated.

\textbf{Definition 2 (Jacobian-Vector Product).}  Given a differentiable
function $\mathbf{f}: \mathbb{R}^n \to \mathbb{R}^m$ evaluated at a point
$\mathbf{x}_0 \in \mathbb{R}^n$, and a tangent vector
$\mathbf{v} \in \mathbb{R}^n$, the Jacobian-vector product (JVP) is:
\begin{equation}
  \jvp\bigl(\mathbf{f};\;\mathbf{x}_0;\;\mathbf{v}\bigr)
    \coloneqq J_\mathbf{f}(\mathbf{x}_0)\,\mathbf{v}
    = \left.\frac{d}{d\alpha}\,\mathbf{f}(\mathbf{x}_0 + \alpha\,\mathbf{v})
      \right|_{\alpha=0}
    \;\in \mathbb{R}^m
  \label{eq:jvp_operator}
\end{equation}
where $J_\mathbf{f}(\mathbf{x}_0) \in \mathbb{R}^{m \times n}$ is the
Jacobian matrix.  The JVP can be computed in $O(n)$ time via forward-mode
automatic differentiation (\code{torch.func.jvp}), without ever constructing
the full $O(n^2)$ Jacobian.

In our setting, $\mathbf{f} = u_\theta$ maps
$(\underbrace{z_t}_{\mathbb{R}^d},\;\underbrace{t}_{\mathbb{R}},\;
\underbrace{r}_{\mathbb{R}})$ to $\mathbb{R}^d$ with $d = 442{,}368$, so the
Jacobian $J \in \mathbb{R}^{442{,}368 \times 442{,}370}$ is never formed.
The tangent vector is $(\tilde{v},\,1,\,0)$ (see \S\ref{sec:compound_v}).

%───────────────────────────────────────────────────────────────────────────────
\subsection{Flow Matching Preliminaries}

Flow matching~\cite{lipman2023,liu2023} defines a probability path between data
and noise via linear interpolation.  For a single sample:
\begin{equation}
  z_t = (1-t)\,z_0 + t\,\varepsilon, \qquad
  z_0 \sim p_{\mathrm{data}},\quad
  \varepsilon \sim \mathcal{N}(\mathbf{0},\,I_d),\quad
  t \in [0,1]
  \label{eq:interpolation}
\end{equation}
\begin{shapes}
$z_0, \varepsilon, z_t \in \mathbb{R}^{C \times D \times H \times W}$;
\quad batched: $(B, 4, 48, 48, 48)$.
\quad $I_d$ is the $d \times d$ identity with $d = 442{,}368$.
\quad $t \in \mathbb{R}$; batched: $(B,)$, broadcast to $(B,1,1,1,1)$.\\
\textbf{Code:} \code{meanflow\_loss.py:201}:
\code{z\_t = (1 - t\_broad) * z\_0 + t\_broad * eps}
where \code{t\_broad = t.view(-1, 1, 1, 1, 1)}.
\end{shapes}

\noindent
where $t{=}0$ is data and $t{=}1$ is noise. The \textbf{conditional velocity field}
is defined as:
\begin{equation}
  v_c(z_t,t) \coloneqq \varepsilon - z_0
  \quad \in \mathbb{R}^{C \times D \times H \times W}
  \label{eq:cond_velocity}
\end{equation}
\begin{shapes}
$v_c \in \mathbb{R}^{C \times D \times H \times W}$;
\quad batched: $(B, 4, 48, 48, 48)$.\\
\textbf{Code:} \code{meanflow\_loss.py:204}: \code{v\_c = eps - z\_0}.
\end{shapes}

\noindent
A neural network $v_\theta$ is trained to match this field:
\begin{equation}
  \mathcal{L}_{\mathrm{FM}}
    = \mathbb{E}_{z_0,\,\varepsilon,\,t}
      \left[\,\bigl\|v_\theta(z_t,t) - v_c\bigr\|_2^2\,\right]
  \label{eq:fm_loss}
\end{equation}
where $\|\cdot\|_2^2$ denotes the squared $L_2$ norm summed over all
$d = C \cdot D \cdot H \cdot W = 442{,}368$ elements.

\textbf{Limitation.}  Sampling requires integrating the learned velocity field
from $t{=}1$ to $t{=}0$ via numerical ODE solvers, requiring $K \geq 5$
network evaluations even with trajectory straightening (rectified flow).

%───────────────────────────────────────────────────────────────────────────────
\subsection{MeanFlow: Average Velocity for 1-Step Generation}

MeanFlow~\cite{meanflow} replaces the instantaneous velocity $v(z_t,t)$ with
the \textbf{average velocity} over an interval $[r,t]$:
\begin{equation}
  u(z_t,r,t)
    \coloneqq \frac{1}{t-r}\int_r^t v(z_s,s)\,ds
    \quad \in \mathbb{R}^{C \times D \times H \times W}
  \label{eq:avg_velocity}
\end{equation}
\begin{shapes}
$u \in \mathbb{R}^{C \times D \times H \times W}$;
\quad batched: $(B, 4, 48, 48, 48)$.
\quad $r, t \in \mathbb{R}$; batched: $(B,)$.
\end{shapes}

\noindent
If the average velocity from $t{=}0$ to $t{=}1$ is known exactly, the entire
flow can be computed in one step:
\begin{equation}
  z_0 = z_1 - u_\theta(z_1,\,r{=}0,\,t{=}1)
  \quad \in \mathbb{R}^{C \times D \times H \times W}
  \qquad \text{[1-NFE generation]}
  \label{eq:one_nfe}
\end{equation}
\begin{shapes}
$z_1 = \varepsilon \sim \mathcal{N}(\mathbf{0}, I_d)$: $(B, 4, 48, 48, 48)$.
\quad $r = \mathbf{0}_B$, $t = \mathbf{1}_B$: each $(B,)$.\\
\textbf{Code:} \code{one\_step.py:37}: \code{output = model(noise, r, t)}.
\end{shapes}

%───────────────────────────────────────────────────────────────────────────────
\subsection{The MeanFlow Identity and Compound Velocity}
\label{sec:compound_v}

The average velocity $u$ must satisfy a self-consistency condition.
Differentiating the integral definition~\eqref{eq:avg_velocity} with respect
to $t$ and applying the chain rule yields the \textbf{MeanFlow Identity}:
\begin{equation}
  v(z_t,t)
    = u(z_t,r,t)
      + (t-r)\!\left[
          \underbrace{\frac{\partial u}{\partial z_t}\,v(z_t,t)}_{
            \in\,\mathbb{R}^{d},\;\text{spatial JVP}}
          + \underbrace{\frac{\partial u}{\partial t}}_{
            \in\,\mathbb{R}^{d},\;\text{temporal deriv.}}
        \right]
  \label{eq:mf_identity}
\end{equation}
\begin{shapes}
All terms $\in \mathbb{R}^{C \times D \times H \times W}$; batched: $(B, 4, 48, 48, 48)$.\\
$\frac{\partial u}{\partial z_t} \in \mathbb{R}^{d \times d}$ is the Jacobian
(never formed); only its product with $v$ is computed via JVP.\\
$(t-r) \in \mathbb{R}$; batched: $(B,)$, broadcast to $(B,1,1,1,1)$.
\end{shapes}

\noindent
The right-hand side, evaluated with the neural approximation $u_\theta$, gives
the \textbf{compound velocity}.  We first compute the directional derivative
(JVP) of $u_\theta$ with respect to its inputs $(z_t, t, r)$ in the tangent
direction $(\tilde{v}, 1, 0)$:
\begin{equation}
  \frac{du_\theta}{dt}\bigg|_{\tilde{v}}
    \coloneqq \jvp\!\bigl(u_\theta;\;(z_t,t,r);\;(\tilde{v},1,0)\bigr)
    = \frac{\partial u_\theta}{\partial z_t}\,\tilde{v}
      + \frac{\partial u_\theta}{\partial t}\cdot 1
      + \frac{\partial u_\theta}{\partial r}\cdot 0
    \;\in \mathbb{R}^{C \times D \times H \times W}
  \label{eq:jvp_def}
\end{equation}
\begin{shapes}
Primals: $z_t \in \mathbb{R}^{B \times 4 \times 48 \times 48 \times 48}$,
$t \in \mathbb{R}^{B}$, $r \in \mathbb{R}^{B}$.\\
Tangents: $\tilde{v} \in \mathbb{R}^{B \times 4 \times 48 \times 48 \times 48}$,
$\mathbf{1}_B \in \mathbb{R}^{B}$, $\mathbf{0}_B \in \mathbb{R}^{B}$.\\
Output: $\frac{du}{dt} \in \mathbb{R}^{B \times 4 \times 48 \times 48 \times 48}$.\\
\textbf{Code:} \code{jvp\_strategies.py:117}: \code{u, du\_dt, v = torch.func.jvp(}\\
\code{\quad \_u\_with\_v\_aux, (z\_t, t, r), (v\_tangent, dt, dr), has\_aux=True)}.\\
Here \code{dt = torch.ones\_like(t)}: $(B,)$; \code{dr = torch.zeros\_like(r)}: $(B,)$.
\end{shapes}

\noindent
The compound velocity is then:
\begin{equation}
  \boxed{
  V_\theta
    \coloneqq u_\theta
      + (t - r) \cdot \sg\!\left[\frac{du_\theta}{dt}\bigg|_{\tilde{v}}\right]
  }
  \quad \in \mathbb{R}^{C \times D \times H \times W}
  \label{eq:compound_V}
\end{equation}
\begin{shapes}
$u_\theta, V_\theta \in \mathbb{R}^{B \times 4 \times 48 \times 48 \times 48}$;
$\sg[\cdot]$ is \code{.detach()} (Def.~1, Eq.~\ref{eq:sg_def}).\\
$(t-r)$: $(B,)$ broadcast to $(B,1,1,1,1)$ via \code{.view(-1,1,1,1,1)}.\\
\textbf{Code:} \code{jvp\_strategies.py:59}:
\code{V = u + t\_minus\_r * du\_dt.detach()}.\\
Gradients flow through $u_\theta$ but \textbf{not} through $du/dt$ (stop-gradient).
\end{shapes}

\noindent
The $\sg[\cdot]$ on the JVP term is essential: it ensures that during
backpropagation, gradients flow only through $u_\theta$ (left term), not
through the JVP correction.  Without $\sg[\cdot]$, second-order derivatives of
$u_\theta$ would be needed, making training intractable.

\textbf{Tangent vector $\tilde{v}$.}  In the improved MeanFlow (iMF)
formulation~\cite{imf}, the tangent $\tilde{v}$ is the model's own
instantaneous velocity estimate at $r{=}t$ (see \S\ref{sec:vhead} for
the dual-head variant):
\begin{equation}
  \tilde{v} = u_\theta(z_t,\,r{=}t,\,t)
  \quad \in \mathbb{R}^{C \times D \times H \times W}
  \label{eq:tangent_imf}
\end{equation}
This makes $V_\theta$ depend only on $z_t$ (not on ground-truth $z_0$ or
$\varepsilon$), matching the inference setting where data is unavailable.

\textbf{Special cases.}  When $r = t$ (flow-matching samples, 50\% of the
batch), $(t-r) = 0$ and $V_\theta = u_\theta$, recovering standard flow
matching.  When $r < t$ (MeanFlow samples), the JVP correction enforces
self-consistency of the average velocity across intervals.

%───────────────────────────────────────────────────────────────────────────────
\subsection{Improved MeanFlow (iMF) and the Dual-Head Architecture}
\label{sec:vhead}

The original MeanFlow uses the ground-truth velocity
$v_c = \varepsilon - z_0$ as the JVP tangent.  \textbf{Problem:} this creates
a dependency on data at inference time (where $z_0$ is unavailable).  The
improved MeanFlow (iMF)~\cite{imf} resolves this by using the model's own
prediction as the tangent, making $V$ a function of $z_t$ alone.

\textbf{The dual-head extension} (our architecture, inspired by iMF) introduces
two output heads on a shared UNet backbone:
\begin{itemize}[nosep]
  \item \textbf{u-head} (main): Predicts the denoised data estimate $\hat{x}$
    (x-prediction mode), from which $u_\theta$ is derived.  Used at inference.
  \item \textbf{v-head} (auxiliary): Predicts the instantaneous velocity
    $\hat{v}$, directly supervised against $v_c$.  Provides the JVP tangent
    $\tilde{v} = \hat{v}$ during training.  \textbf{Disabled at inference}
    (zero cost).
\end{itemize}

The shared backbone produces a feature tensor $\mathbf{h} \in
\mathbb{R}^{B \times 64 \times 48 \times 48 \times 48}$ (channel dimension
equals \code{channels[0]}$= 64$).  Both heads branch from $\mathbf{h}$:
\begin{equation}
  \hat{x} = \mathrm{UNet.out}(\mathbf{h})
    \;\in \mathbb{R}^{C \times D \times H \times W},
  \qquad
  \hat{v} = \mathrm{v\_out}(\mathbf{h})
    \;\in \mathbb{R}^{C \times D \times H \times W}
  \label{eq:dual_heads}
\end{equation}
\begin{shapes}
$\mathbf{h}$: $(B, 64, 48, 48, 48)$ --- shared backbone final feature map.\\
$\hat{x}$: $(B, 4, 48, 48, 48)$ via \code{self.unet.out(h)} (Conv3d $64 \to 4$).\\
$\hat{v}$: $(B, 4, 48, 48, 48)$ via \code{self.v\_out(h)} (ResBlock $\to$ GN $\to$ SiLU $\to$ Conv3d $64 \to 4$).\\
\textbf{Code:} \code{maisi\_unet.py:415--418}.
\end{shapes}

\noindent
The v-head solves a critical practical problem: without it, the tangent comes
from the u-head's own prediction, which is poor in early training, creating a
bootstrapping problem that causes loss divergence.  The v-head receives direct
supervision ($\|\hat{v} - v_c\|^p$), providing high-quality tangents from the
first epoch.  In our training, the v-head cosine alignment $\cos(\hat{v},v_c)$
reached 0.39 by epoch~50---85\% of its final value of~0.44.

%───────────────────────────────────────────────────────────────────────────────
\subsection{x-Prediction Reparameterisation}
\label{sec:xpred}

Instead of directly outputting the average velocity $u$, the u-head outputs a
denoised data estimate $\hat{x}$ (x-prediction, following pMF~\cite{pmf}).
The average velocity is derived via:
\begin{equation}
  \hat{x} = f_\theta(z_t,r,t)
    \;\in \mathbb{R}^{C \times D \times H \times W},
  \qquad
  u_\theta = \frac{z_t - \hat{x}}{\max(t,\,t_{\min})}
    \;\in \mathbb{R}^{C \times D \times H \times W}
  \label{eq:xpred}
\end{equation}
where $t_{\min} = 0.05$ is a clamping threshold that prevents the $1/t$
singularity near $t = 0$.  The division is element-wise after broadcasting
$t$ to shape $(B, 1, 1, 1, 1)$.
\begin{shapes}
$\hat{x}, z_t, u_\theta$: each $(B, 4, 48, 48, 48)$.
\quad $t$: $(B,)$, clamped to $[\,t_{\min},\,\infty)$, broadcast to $(B,1,1,1,1)$.\\
\textbf{Code:} \code{meanflow\_loss.py:113--119}:
\code{t\_safe = t\_.view(-1, *([1]*(z.ndim-1))).clamp(min=t\_min)};
\code{return (z - x\_hat) / t\_safe}.
\end{shapes}

\textbf{Justification (manifold hypothesis).}  The x-prediction target lies on
the data manifold, which has low intrinsic dimensionality.  The velocity $u$
spans a higher-dimensional space.  For architectures with a bottleneck (UNet
encoder-decoder), predicting a low-dimensional target ($\hat{x}$) is easier
than predicting a high-dimensional one ($u$).

\textbf{Quantitative criterion} (pMF Table~2): x-prediction dominates when
$d_{\mathrm{input}}/d_{\mathrm{bottleneck}} > 1$.  For our 3D UNet:
\begin{equation}
  \frac{d_{\mathrm{input}}}{d_{\mathrm{bottleneck}}}
    = \frac{C \times D \times H \times W}{\text{channels}[3] \times (D/8)^3}
    = \frac{4 \times 48^3}{512 \times 6^3}
    = \frac{442{,}368}{110{,}592} \approx 4
\end{equation}
This is firmly in the x-prediction regime.

\textbf{At inference (1-NFE).}  With x-prediction, 1-step sampling simplifies
to a single forward pass---the model directly outputs $\hat{x} = z_0$:
\begin{equation}
  z_0 = f_\theta(\varepsilon,\;r{=}0,\;t{=}1)
    \;\in \mathbb{R}^{C \times D \times H \times W}
  \label{eq:1nfe_xpred}
\end{equation}
\begin{shapes}
$\varepsilon$: $(B, 4, 48, 48, 48)$.
\quad $r = \mathbf{0}_B$: $(B,)$; $t = \mathbf{1}_B$: $(B,)$.
\quad Output: $(B, 4, 48, 48, 48)$.\\
\textbf{Code:} \code{one\_step.py:37--44}:
\code{output = model(noise, r, t)}; \code{return output} (x-pred: direct).
\end{shapes}

%───────────────────────────────────────────────────────────────────────────────
\subsection{Latent Space Formulation}

All of the above operates identically in latent space.  The computational
advantages are:
\begin{enumerate}[nosep]
  \item \textbf{JVP cost reduction:} Latent dimension
    $d = 4\times 48^3 = 442{,}368$ vs pixel space
    $d = 192^3 = 7{,}077{,}888$---a \textbf{16$\times$ reduction} in JVP
    compute per iteration.
  \item \textbf{Memory reduction:} The UNet operates on $48^3$ feature maps
    vs $192^3$---approximately \textbf{64$\times$ reduction} in activation
    memory per JVP pass.
  \item \textbf{Training efficiency:} The latent space is pre-computed
    (Phase~1), so VAE encode/decode cost is amortised across all training
    epochs.
\end{enumerate}


%═══════════════════════════════════════════════════════════════════════════════
\section{Architecture Design}
\label{sec:architecture}

\subsection{Frozen MAISI VAE (Encoder-Decoder)}

The MAISI VAE~\cite{maisi} is a 3D variational autoencoder with adversarial
training, pre-trained on ${\sim}55$K CT+MRI volumes.  We use it as a
\textbf{frozen foundation model}---all parameters are locked, and it is never
fine-tuned. Key properties are listed in Table~\ref{tab:vae}.

\begin{table}[ht]
\centering
\caption{MAISI VAE properties.}
\label{tab:vae}
\begin{tabular}{ll}
\toprule
\textbf{Property}      & \textbf{Value} \\
\midrule
Parameters             & 20,944,897 (${\sim}$21\,M) \\
Encoder stages         & 3 levels of $2\times$ strided 3D conv \\
Latent channels        & 4 (with KL regularisation) \\
Spatial compression    & $4\times$ per axis: $(1,192^3)\!\to\!(4,48^3)$ \\
Attention              & None (all \code{attention\_levels=false}) \\
Training losses        & $L_1$ + LPIPS + PatchGAN + KL \\
Checkpoint format      & Wrapped in \code{"unet\_state\_dict"} key \\
Scale factor           & $s = 0.96240234375$ \\
Memory optimisation    & \code{num\_splits} for chunk-based processing \\
\bottomrule
\end{tabular}
\end{table}

The VAE encoding and decoding operations are:
\begin{equation}
  z = s \cdot \mathrm{Enc}_\phi(x_\mathrm{MRI})
    \;\in \mathbb{R}^{4 \times 48 \times 48 \times 48},
  \qquad
  \hat{x}_\mathrm{MRI} = \mathrm{Dec}_\phi(z\,/\,s)
    \;\in \mathbb{R}^{1 \times 192 \times 192 \times 192}
  \label{eq:vae}
\end{equation}
\begin{shapes}
$x_\mathrm{MRI}$: $(B, 1, 192, 192, 192)$ --- single-channel T1 MRI.\\
$z$: $(B, 4, 48, 48, 48)$ --- latent ($4\times$ spatial compression per axis, $1 \to 4$ channels).\\
$s = 0.9624$ --- scale factor extracted from
\code{diff\_unet\_3d\_rflow-mr.pt["scale\_factor"]}.
\end{shapes}

\textbf{Reconstruction quality} (Phase~0, 20 IXI volumes): Mean SSIM\,=\,0.9213,
Mean PSNR\,=\,30.86\,dB.

\textbf{Latent distribution quality.}  Encoding all 6,471~volumes (train + val
+ test) confirms the latent space is well-regularised
(Table~\ref{tab:latent_stats}).

\begin{table}[ht]
\centering
\caption{Per-channel latent statistics ($n = 6{,}471$ encoded volumes).}
\label{tab:latent_stats}
\begin{tabular}{ccccc}
\toprule
\textbf{Channel} & \textbf{Mean} & \textbf{Std} & \textbf{Skewness} & \textbf{Kurtosis} \\
\midrule
0 & $-0.053$ & $0.970$ & $+0.105$ & $-0.123$ \\
1 & $-0.185$ & $1.019$ & $+0.002$ & $-0.019$ \\
2 & $-0.051$ & $0.970$ & $+0.045$ & $+0.099$ \\
3 & $+0.001$ & $1.011$ & $+0.069$ & $+0.144$ \\
\bottomrule
\end{tabular}
\end{table}

The distributions are near-Gaussian: skewness $|\gamma| < 0.11$, excess
kurtosis $|\kappa| < 0.15$, and standard deviations cluster around
$0.97$--$1.02$.  Cross-channel correlations are negligible
(max off-diagonal $|r| = 0.046$), confirming the KL regularisation produces
approximately independent, unit-variance channels.

\subsection{MeanFlow UNet (Generative Model)}

The MeanFlow UNet (\code{MAISIUNetWrapper}) uses the \textbf{same architecture}
as the MAISI diffusion UNet but with random initialisation and custom dual-time
conditioning.  Key properties are listed in Table~\ref{tab:unet}.

\begin{table}[ht]
\centering
\caption{MeanFlow UNet properties and tensor shapes through the encoder-decoder path.}
\label{tab:unet}
\begin{tabular}{lll}
\toprule
\textbf{Property}          & \textbf{Value} & \textbf{Feature map shape} \\
\midrule
Backbone                   & MONAI \code{DiffusionModelUNet} (3D) & --- \\
Total parameters           & ${\sim}$178\,M & --- \\
\code{conv\_in}            & Conv3d($4, 64, 3$, pad$=1$) & $(B, 64, 48, 48, 48)$ \\
Down block 0 (ch$=64$)     & 2 ResBlocks, no attention & $(B, 64, 24, 24, 24)$ \\
Down block 1 (ch$=128$)    & 2 ResBlocks, no attention & $(B, 128, 12, 12, 12)$ \\
Down block 2 (ch$=256$)    & 2 ResBlocks + Transformer & $(B, 256, 6, 6, 6)$ \\
Down block 3 (ch$=512$)    & 2 ResBlocks + Transformer & $(B, 512, 3, 3, 3)$ \\
Middle block (ch$=512$)    & ResBlock + Attn + ResBlock & $(B, 512, 3, 3, 3)$ \\
Up blocks (mirror down)    & With skip connections & up to $(B, 64, 48, 48, 48)$ \\
\code{unet.out} (u-head)   & GN + SiLU + Conv3d($64, 4$) & $(B, 4, 48, 48, 48)$ \\
\code{v\_out} (v-head)     & ResBlock + GN + SiLU + Conv3d($64, 4$) & $(B, 4, 48, 48, 48)$ \\
\midrule
Attention heads            & 32 channels/head at levels 2--3 & --- \\
GroupNorm groups           & 32 & --- \\
Flash attention            & \textbf{Disabled} (\code{torch.func.jvp} compat.) & --- \\
Grad.\ checkpointing      & \textbf{Disabled} (exact JVP compat.) & --- \\
Prediction type            & x-prediction ($\hat{x}$) & --- \\
\bottomrule
\end{tabular}
\end{table}

\textbf{Flash attention incompatibility.}  \code{torch.func.jvp} uses
forward-mode AD, which requires the computation graph to be fully traceable.
Flash attention's fused CUDA kernels are opaque to PyTorch's AD system.
Disabling flash attention adds ${\sim}10\%$ latency but enables exact JVP.

\textbf{Gradient checkpointing incompatibility.}  Gradient checkpointing
re-executes forward passes during backward, which conflicts with the
forward-mode AD tape maintained by \code{torch.func.jvp}.

\textbf{In-place operation patching.}  \code{torch.func.jvp} cannot handle
in-place tensor mutations.  The function \code{patch\_inplace\_ops()} (line~133
of \code{maisi\_unet.py}) recursively replaces \code{nn.SiLU(inplace=True)}
and \code{nn.ReLU(inplace=True)} with out-of-place variants throughout the
model.

\subsection{Dual-Time Conditioning}

MeanFlow requires conditioning on two time variables: $t$ (current time) and
$r$ (interval lower bound).  We implement three conditioning modes and select
\code{t\_h} based on ablation (Table~\ref{tab:conditioning}).

\begin{table}[ht]
\centering
\caption{Conditioning modes.  Selected mode in \textbf{bold}.}
\label{tab:conditioning}
\begin{tabular}{llll}
\toprule
\textbf{Mode} & \textbf{Inputs} & \textbf{Embedding} & \textbf{Source} \\
\midrule
\code{dual} & $(r,t)$ & $\mathbf{e}_r+\mathbf{e}_t$ via separate MLPs & pMF \\
\code{h}    & $h=t{-}r$ & $\mathbf{e}_h$ via UNet MLP & iMF \\
\textbf{\code{t\_h}} & \textbf{$(t,\,h{=}t{-}r)$} & \textbf{$\mathbf{e}_t + \mathbf{e}_h$ via two MLPs} & \textbf{MF Tab.\,1c} \\
\bottomrule
\end{tabular}
\end{table}

The \code{t\_h} mode conditions on both the absolute time $t$ and the interval
width $h = t-r$.  The original MeanFlow paper (Table~1c) reports FID~61.06
for $(t,h)$ vs~63.13 for $h$-only on ImageNet.

\textbf{Sinusoidal embedding.}  Given a scalar $\tau \in [0,1]$ (either $t$ or
$h$), the sinusoidal positional embedding~\cite{ddpm} is defined as:
\begin{equation}
  \mathrm{SinEmb}(\tau)
    \coloneqq \bigl[\cos(\omega_j \cdot \tau'),\;\sin(\omega_j \cdot \tau')\bigr]_{j=0}^{d_\mathrm{sin}/2-1}
    \;\in \mathbb{R}^{d_\mathrm{sin}}
  \label{eq:sin_emb}
\end{equation}
where $\tau' = 1000 \cdot \tau$ (time scaling), $d_\mathrm{sin} = 64$ (=
\code{channels[0]}), and the frequencies are
$\omega_j = \exp(-\log(10000) \cdot j / (d_\mathrm{sin}/2))$ for
$j = 0, \ldots, 31$.
\begin{shapes}
$\mathrm{SinEmb}(\tau)$: $(B, 64)$.\\
\textbf{Code:} \code{maisi\_unet.py:48--58},
\code{get\_timestep\_embedding(timesteps * 1000, embedding\_dim=64)}.
\end{shapes}

\textbf{Time embedding MLP.}  Each sinusoidal embedding is projected through a
2-layer MLP:
\begin{equation}
  \mathrm{MLP}: \mathbb{R}^{d_\mathrm{sin}} \to \mathbb{R}^{d_\mathrm{emb}},
  \qquad
  \mathbf{e} = W_2\,\mathrm{SiLU}(W_1\,\mathrm{SinEmb}(\tau) + \mathbf{b}_1) + \mathbf{b}_2
  \label{eq:time_mlp}
\end{equation}
where $W_1 \in \mathbb{R}^{d_\mathrm{emb} \times d_\mathrm{sin}}$,
$W_2 \in \mathbb{R}^{d_\mathrm{emb} \times d_\mathrm{emb}}$,
$d_\mathrm{sin} = 64$, $d_\mathrm{emb} = 256$.

The combined conditioning embedding (in \code{t\_h} mode) is:
\begin{equation}
  \mathbf{e} = \mathrm{MLP}_t\!\bigl(\mathrm{SinEmb}(t)\bigr)
             + \mathrm{MLP}_h\!\bigl(\mathrm{SinEmb}(h)\bigr)
    \;\in \mathbb{R}^{d_\mathrm{emb}}
  \label{eq:conditioning}
\end{equation}
\begin{shapes}
$\mathrm{SinEmb}(t)$, $\mathrm{SinEmb}(h)$: each $(B, 64)$.\\
$\mathrm{MLP}_t$ output, $\mathrm{MLP}_h$ output: each $(B, 256)$.\\
$\mathbf{e}$: $(B, 256)$ --- injected into every ResBlock and Transformer block.\\
\textbf{Code:} \code{maisi\_unet.py:448--460}:
\code{emb = t\_emb + h\_emb}.\\
$\mathrm{MLP}_t$: \code{self.unet.time\_embed} (MONAI built-in);
$\mathrm{MLP}_h$: \code{self.h\_embed} (custom).
\end{shapes}

\subsection{v-Head (Auxiliary Tangent Predictor)}
\label{sec:vhead_arch}

The v-head is a lightweight auxiliary output path branching from the UNet's
final feature map $\mathbf{h}$ (before the main output convolution):
\begin{center}
\begin{BVerbatim}
Shared Backbone (B, 64, 48, 48, 48)
  |-- u-head: UNet out conv  --> (B, 4, 48, 48, 48)  [main]
  |-- v-head: ResBlock -> GN -> SiLU -> Conv3d --> (B, 4, 48, 48, 48)
\end{BVerbatim}
\end{center}

The v-head ResBlock~(\code{\_VHeadResBlock}, \code{maisi\_unet.py:205--235})
has the following structure:
\begin{equation}
  \mathrm{ResBlock}(\mathbf{h})
    = \mathbf{h}
      + \underbrace{
          \mathrm{Conv}_{3 \times 3 \times 3}^{64 \to 64}\bigl(
            \mathrm{SiLU}\bigl(
              \mathrm{GN}_{32}\bigl(
                \mathrm{Conv}_{3 \times 3 \times 3}^{64 \to 64}\bigl(
                  \mathrm{SiLU}\bigl(\mathrm{GN}_{32}(\mathbf{h})\bigr)
                \bigr)
              \bigr)
            \bigr)
          \bigr)
        }_{\text{zero-initialised second conv}}
  \label{eq:vhead_resblock}
\end{equation}
\begin{shapes}
$\mathbf{h} \in \mathbb{R}^{B \times 64 \times 48 \times 48 \times 48}$.
\quad $\mathrm{GN}_{32}$: GroupNorm with 32 groups over 64 channels.\\
Conv weights: $\mathbb{R}^{64 \times 64 \times 3 \times 3 \times 3}$
($= 110{,}592$ params each). Second conv zero-initialised.\\
Total v-head: ${\sim}228$\,K params ($<0.13\%$ of 178\,M backbone).
\end{shapes}

\begin{table}[ht]
\centering
\caption{v-Head properties.}
\label{tab:vhead}
\begin{tabular}{ll}
\toprule
\textbf{Property}   & \textbf{Value} \\
\midrule
ResBlocks            & 1 \\
Parameters           & ${\sim}$228\,K ($<$0.13\% of total 178\,M) \\
Initialisation       & Final Conv3d zero-initialised (weights and bias) \\
Inference cost       & Zero (output discarded at sampling) \\
\bottomrule
\end{tabular}
\end{table}

\textbf{Zero initialisation rationale.}  The v-head starts at identically zero
output ($\hat{v} = \mathbf{0}$), so the u-head gradient signal is unperturbed
at initialisation.  Our training confirms:
$\cos(\hat{v},v_c)$ rises from 0.17 (epoch~0) to 0.39 (epoch~50) to 0.44
(epoch~689).


%═══════════════════════════════════════════════════════════════════════════════
\section{Loss Function and Training Objective}
\label{sec:loss}

\subsection{Per-Channel \texorpdfstring{$L_p$}{Lp} Loss}
\label{sec:lp_loss}

The base loss function computes a per-channel, spatially-summed $L_p$ norm.
For tensors $\hat{y}, y \in \mathbb{R}^{C \times D \times H \times W}$:
\begin{equation}
  \ell_p(\hat{y},\,y)
    \coloneqq \sum_{c=1}^{C}\;\sum_{i=1}^{D}\sum_{j=1}^{H}\sum_{k=1}^{W}
      \bigl|\hat{y}_{c,i,j,k} - y_{c,i,j,k}\bigr|^p
    \;\in \mathbb{R}_{\ge 0}
  \label{eq:lp_loss}
\end{equation}
\begin{shapes}
$\hat{y}, y \in \mathbb{R}^{B \times 4 \times 48 \times 48 \times 48}$.\\
Output with \code{reduction="none"}: per-sample scalar $(B,)$---summed over all
$C \times D \times H \times W = 4 \times 48^3 = 442{,}368$ elements per sample.\\
\textbf{Code:} \code{lp\_loss.py:34--48}:
\code{diff = (pred - target).abs()}; \code{diff = diff.pow(p)};
\code{per\_sample = diff.sum(dim=[1,2,3,4])}.
\end{shapes}

\noindent
For the best model: $p{=}2.0$ (standard $L_2$), no per-channel weights.
When optional channel weights $\mathbf{w} \in \mathbb{R}^C$ are provided:
\begin{equation}
  \ell_p^{\mathbf{w}}(\hat{y},y)
    = \sum_{c=1}^{C} w_c \!\sum_{i,j,k}
      |\hat{y}_{c,i,j,k} - y_{c,i,j,k}|^p
  \label{eq:lp_weighted}
\end{equation}
This extends the SLIM-Diff per-channel loss~\cite{slimdiff} from pixel-space
DDPM to latent-space MeanFlow.

\subsection{Combined iMF Dual-Head Loss}
\label{sec:dual_loss}

The training loss has two independently-weighted components:
\begin{equation}
  \boxed{
  \mathcal{L}
    = \frac{1}{B}\sum_{b=1}^{B}\!\left[
        \frac{\ell_p^{(b)}(V_\theta,\,v_c)}{w_u^{(b)}}
        + \frac{\ell_p^{(b)}(\hat{v},\,v_c)}{w_v^{(b)}}
      \right]
  }
  \;\in \mathbb{R}
  \label{eq:total_loss}
\end{equation}
where the superscript $(b)$ denotes the $b$-th sample in the batch, and:
\begin{itemize}[nosep]
  \item $\ell_p^{(b)}(V_\theta, v_c)$ is the \textbf{compound velocity loss}
    (u-head): enforces MeanFlow self-consistency.
  \item $\ell_p^{(b)}(\hat{v}, v_c)$ is the \textbf{tangent supervision loss}
    (v-head): directly supervises the auxiliary head.
  \item $w_u^{(b)}, w_v^{(b)}$ are per-sample adaptive weights (defined below).
\end{itemize}
Both losses target the same conditional velocity $v_c = \varepsilon - z_0
\in \mathbb{R}^{C \times D \times H \times W}$.
\begin{shapes}
$V_\theta, \hat{v}, v_c$: each $(B, 4, 48, 48, 48)$.\\
$\ell_p^{(b)}(\cdot, \cdot)$: scalar per sample, so the full vector is $(B,)$.\\
$\mathcal{L}$: scalar (mean over batch).\\
\textbf{Code:} \code{meanflow\_loss.py:342--356}:
\code{raw\_loss\_u = lp\_loss(V, v\_c, ...)};
\code{raw\_loss\_v = lp\_loss(v, v\_c, ...)};
\code{loss = (loss\_u + loss\_v).mean()}.
\end{shapes}

\subsection{Adaptive Weighting}
\label{sec:adaptive}

Each loss component is independently normalised by its own magnitude to prevent
either head from dominating the gradient.  For a generic per-sample raw loss
$\ell^{(b)} \in \mathbb{R}_{\ge 0}$, the adaptive weight and weighted loss
are:
\begin{equation}
  w^{(b)}
    \coloneqq \bigl(\sg[\ell^{(b)}] + \varepsilon_{\mathrm{norm}}\bigr)^{p_{\mathrm{norm}}}
    \;\in \mathbb{R}_{> 0}
  \label{eq:adaptive_weight}
\end{equation}
\begin{equation}
  \ell^{(b)}_{\mathrm{weighted}}
    \coloneqq \frac{\ell^{(b)}}{w^{(b)}}
    \;\in \mathbb{R}_{\ge 0}
  \label{eq:adaptive_weighted_loss}
\end{equation}
where $\sg[\cdot]$ is the stop-gradient operator (Eq.~\ref{eq:sg_def})
ensuring the weight does not contribute gradients, $\varepsilon_{\mathrm{norm}}
= 1.0$, and $p_{\mathrm{norm}} = 1.0$.  Substituting these defaults:
\begin{equation}
  w^{(b)} = \sg[\ell^{(b)}] + 1.0,
  \qquad
  \ell^{(b)}_{\mathrm{weighted}}
    = \frac{\ell^{(b)}}{\sg[\ell^{(b)}] + 1.0}
  \label{eq:adaptive_simplified}
\end{equation}
\begin{shapes}
$\ell^{(b)}$: per-sample scalar; full batch: $(B,)$.\\
$w^{(b)}$: per-sample scalar; full batch: $(B,)$. Computed with \code{.detach()} on $\ell$.\\
$\ell^{(b)}_{\mathrm{weighted}}$: per-sample scalar; full batch: $(B,)$.\\
\textbf{Code:} \code{meanflow\_loss.py:347--351}:
\code{adp\_weight\_u = (raw\_loss\_u.detach() + norm\_eps) ** norm\_p};
\code{loss\_u = raw\_loss\_u / adp\_weight\_u}.
Applied \textbf{independently} to u-head and v-head.
\end{shapes}

\noindent
This is a form of \textbf{logarithmic loss normalisation}: the effective loss
behaves as $\ell_\mathrm{w} \approx \ell/(\ell+1)$, which is bounded in
$[0,1)$ and varies slowly with raw loss magnitude:
\begin{equation}
  \lim_{\ell \to 0}\;\frac{\ell}{\ell + 1} = 0,
  \qquad
  \lim_{\ell \to \infty}\;\frac{\ell}{\ell + 1} = 1
\end{equation}

At convergence (epoch~388, best FID), the loss decomposition shows this
balancing in action (Table~\ref{tab:loss_decomp}).

\begin{table}[ht]
\centering
\caption{Loss decomposition at best FID (epoch~388).}
\label{tab:loss_decomp}
\begin{tabular}{lrr}
\toprule
\textbf{Component} & \textbf{Raw Loss} & \textbf{Share} \\
\midrule
FM loss ($r=t$, flow matching)       & 715,432   & 10.3\% \\
MF loss ($r<t$, self-consistency)    & 6,254,997 & 89.7\% \\
\midrule
v-head loss                          & 714,646   & --- \\
u-head loss (compound $V$)           & 3,485,215 & --- \\
\bottomrule
\end{tabular}
\end{table}

The MF loss is ${\sim}9\times$ larger than the FM loss, which is expected: the
compound velocity $V$ must also capture the JVP correction term, a harder
target than the direct velocity at $r{=}t$.  The adaptive weighting equalises
their gradient contributions to ${\sim}1.0$ each.

\textbf{Why $\varepsilon_{\mathrm{norm}}{=}1.0$?}  Through Phase~4c debugging,
we discovered that $\varepsilon_{\mathrm{norm}}{=}0.01$ caused catastrophic
gradient amplification for samples with small raw loss.  The weight
$1/(\ell+0.01)$ can reach $100\times$, creating a positive feedback loop.

\textbf{Why $p_{\mathrm{norm}}{=}1.0$?}  With $p_{\mathrm{norm}}{=}0.5$, the
weight becomes $w = \sqrt{\ell+1.0}$, which under-normalises large
losses, leading to a $1000\times$ gradient explosion (Phase~4e).

\subsection{JVP Strategies}
\label{sec:jvp_strat}

Two strategies are implemented for computing the directional derivative in
Eq.~\eqref{eq:jvp_def}.

\textbf{Strategy 1: Exact JVP} (\code{torch.func.jvp}).  Computes the exact
directional derivative via forward-mode automatic differentiation:
\begin{equation}
  \bigl(u_\theta,\;\tfrac{du}{dt},\;\hat{v}\bigr)
    = \texttt{torch.func.jvp}\!\bigl(
        \underbrace{f_{\mathrm{dual}}}_{\mathbb{R}^d \times \mathbb{R} \times \mathbb{R} \,\to\, \mathbb{R}^d \times \mathbb{R}^d},\;
        \underbrace{(z_t,\,t,\,r)}_{\text{primals}},\;
        \underbrace{(\tilde{v},\,\mathbf{1},\,\mathbf{0})}_{\text{tangents}},\;
        \texttt{has\_aux}{=}\texttt{True}\bigr)
  \label{eq:exact_jvp}
\end{equation}
\begin{shapes}
Primals: $z_t \!:\! (B, 4, 48, 48, 48)$; $t \!:\! (B,)$; $r \!:\! (B,)$.\\
Tangents: $\tilde{v} \!:\! (B, 4, 48, 48, 48)$; $\mathbf{1}_B \!:\! (B,)$; $\mathbf{0}_B \!:\! (B,)$.\\
Returns: $u_\theta \!:\! (B, 4, 48, 48, 48)$; $du/dt \!:\! (B, 4, 48, 48, 48)$; $\hat{v} \!:\! (B, 4, 48, 48, 48)$.\\
\code{has\_aux=True} captures $\hat{v}$ without computing its JVP (zero overhead).\\
Compound: $V = u + (t-r) \cdot \sg[du/dt]$: $(B, 4, 48, 48, 48)$.\\
\textbf{Code:} \code{jvp\_strategies.py:109--119}.
Memory: ${\sim}20$\,GB activations per sample at $B{=}2$ on A100-40GB.
\end{shapes}

\noindent
Here $f_{\mathrm{dual}}(z,t,r) \coloneqq (u_\theta(z,t,r),\;\hat{v}(z,t,r))$
returns both the average velocity (through which the JVP is computed) and the
v-head output (captured as auxiliary, no JVP).  The function
\code{\_u\_with\_v\_aux} wraps the dual model output so that
\code{has\_aux=True} passes the v-head through unchanged.

\textbf{Strategy 2: Finite Difference JVP} (FD-JVP).  Approximates the
directional derivative using a perturbation of step size
$\delta = 10^{-3}$:
\begin{equation}
  \frac{du_\theta}{dt}\bigg|_{\tilde{v}}
    \approx \frac{
      u_\theta^{\mathrm{(fp32)}}(z_t + \delta\,\tilde{v},\;t + \delta,\;r)
      - u_\theta^{\mathrm{(fp32)}}(z_t,\;t,\;r)
    }{\delta}
    \;\in \mathbb{R}^{C \times D \times H \times W}
  \label{eq:fd_jvp}
\end{equation}
\begin{shapes}
Perturbed inputs: $z_t + \delta\tilde{v}$: $(B, 4, 48, 48, 48)$; $t + \delta$: $(B,)$.\\
Both $u$ terms cast to fp32 before subtraction (avoids bf16 catastrophic cancellation).\\
Perturbed pass runs under \code{torch.no\_grad()} (since $du/dt$ is always detached).\\
\textbf{Code:} \code{jvp\_strategies.py:156--168}.
\end{shapes}

\textbf{Critical incompatibility (Phase~4f).}  x-prediction + FD-JVP is
\textbf{numerically unstable}.  The x-to-u conversion
$u = (z_t - \hat{x})/t$ has a $1/t$ singularity.  When FD-JVP computes
the difference quotient, the result involves:
\begin{equation}
  \frac{du}{dt}\bigg|_\mathrm{FD}
    = \frac{u(t{+}\delta)-u(t)}{\delta}
    \sim \frac{O(1/t)}{\delta}
    = O\!\left(\frac{1}{t^2\,\delta}\right) \to \infty
  \quad\text{as }t\to t_{\min}
  \label{eq:fd_instability}
\end{equation}
With exact JVP, the $1/t$ factor is analytically differentiated via the
chain rule, yielding stable gradients.

\textbf{Rule:} \fbox{x-pred + exact JVP = stable}\ \ %
\fbox{u-pred + FD-JVP = stable}\ \ %
\fbox{x-pred + FD-JVP = explosion}


%═══════════════════════════════════════════════════════════════════════════════
\section{Data Pipeline}
\label{sec:data}

\subsection{Dataset}

We use 8 datasets from FOMO-60K, a large-scale preprocessed brain MRI
collection (Table~\ref{tab:dataset}).

\begin{table}[ht]
\centering
\caption{Dataset composition from 8 FOMO-60K subsets.}
\label{tab:dataset}
\begin{tabular}{lrrrr}
\toprule
\textbf{Dataset} & \textbf{Train Subj.} & \textbf{Train Scans} & \textbf{Val Scans} & \textbf{Test Scans} \\
\midrule
PT001\_OASIS1 & 352  & 1,414 & 170 & 96 \\
PT002\_OASIS2 & 71   & 682   & 82  & 37 \\
PT005\_IXI    & 494  & 494   & 58  & 29 \\
PT007\_NIMH   & 212  & 428   & 48  & 23 \\
PT008\_DLBS   & 394  & 807   & 109 & 51 \\
PT011\_MBSR   & 125  & 295   & 36  & 16 \\
PT012\_UCLA   & 106  & 106   & 13  & 6  \\
PT015\_NKI    & 722  & 1,245 & 158 & 68 \\
\midrule
\textbf{Total} & \textbf{2,476} & \textbf{5,471} & \textbf{674} & \textbf{326} \\
\bottomrule
\end{tabular}
\end{table}

Many subjects have multiple sessions (longitudinal scans), yielding 5,471
training scans from 2,476 subjects---approximately $4\times$ more data than
the initially planned ${\sim}1{,}172$ scans from 3~datasets.  All volumes are
T1-weighted, skull-stripped, RAS-oriented, and co-registered.  The split is
stratified by dataset with a fixed seed of~42 for reproducibility.

\subsection{Preprocessing Pipeline}

\begin{enumerate}[nosep]
  \item \code{LoadImaged} (MONAI)
  \item \code{EnsureChannelFirstd}
  \item \code{Spacingd}(pixdim=1.0\,mm isotropic, bilinear)
  \item \code{ScaleIntensityRangePercentilesd}(0\%--99.5\% $\to$ [0,\,1])
  \item \code{CropForegroundd}(margin=4) --- brain-centred crop
  \item \code{ResizeWithPadOrCropd}$(192,192,192)$ --- pad to target
  \item \code{EnsureTyped}(dtype=float32)
\end{enumerate}

\textbf{Resolution choice ($192^3$ at 1.0\,mm isotropic).}  Selected after
quantitative analysis of brain extent across 30 FOMO-60K subjects.  Brain AP
extent reaches 193\,mm; $128^3$ and $160^3$ both clip frontal/occipital
cortex.  \code{CropForegroundd} centres the crop on brain tissue, achieving
100\% brain coverage for all subjects.

\textbf{Latent shape:} $192/4 = 48$ per spatial axis, giving latents of shape
$(4,48,48,48)$.

\subsection{Latent Pre-Computation (Phase~1)}

All training volumes are pre-encoded through the frozen MAISI VAE and stored
in HDF5 shards.  A total of 6,471~volumes were encoded (${\sim}$11.5\,GB
storage).

\textbf{Latent statistics} are computed online during encoding via
\code{LatentStatsAccumulator} (sum-of-powers method in float64): per-channel
mean, std, skewness, kurtosis, min, max, and a $4{\times}4$ cross-channel
Pearson correlation matrix.

\subsection{Latent Augmentation}

Three augmentation transforms operate directly on latents during training
(Table~\ref{tab:augmentation}).

\begin{table}[ht]
\centering
\caption{Latent augmentation transforms.}
\label{tab:augmentation}
\begin{tabular}{llll}
\toprule
\textbf{Transform} & \textbf{Prob.} & \textbf{Params} & \textbf{Rationale} \\
\midrule
Flip depth axis   & 0.5 & L--R in RAS        & Bilateral symmetry \\
Gaussian noise    & 0.2 & 5\% of per-ch.\ std & Encoding noise \\
Intensity scale   & 0.2 & $\pm$5\%            & Scanner variation \\
\bottomrule
\end{tabular}
\end{table}

Approximately 68\% of training samples are augmented per epoch
(${\sim}3{,}720$ out of 5,471).


%═══════════════════════════════════════════════════════════════════════════════
\section{Training Protocol}
\label{sec:training}

\subsection{Algorithm (One Training Step)}
\label{sec:algorithm}

Given a mini-batch of $B$ pre-computed latents and the dual-head model
$f_\theta$ (with u-head and v-head), one training step proceeds as follows.
All tensors are annotated with their shapes.

\begin{enumerate}[nosep, leftmargin=2em]
  \item \textbf{Load data:}
    $z_0 \in \mathbb{R}^{B \times 4 \times 48 \times 48 \times 48}$
    (normalised latent from HDF5).

  \item \textbf{Sample noise:}
    $\varepsilon \sim \mathcal{N}(\mathbf{0},\,I_d)
    \in \mathbb{R}^{B \times 4 \times 48 \times 48 \times 48}$.

  \item \textbf{Sample times:}
    $t, r \sim \mathrm{LogitNormal}(\mu{=}{-0.4},\,\sigma{=}1.0)
    \in \mathbb{R}^B$,
    clamped to $[0.001,\,1.0]$, with $t \ge r$.
    For half the batch ($\lfloor B/2 \rfloor$ samples): set $r = t$
    (FM samples).

  \item \textbf{Interpolate:}
    $z_t = (1-\bar{t})\odot z_0 + \bar{t}\odot\varepsilon
    \in \mathbb{R}^{B \times 4 \times 48 \times 48 \times 48}$,
    where $\bar{t} = t.\mathrm{view}(B,1,1,1,1)$.

  \item \textbf{Target:}
    $v_c = \varepsilon - z_0
    \in \mathbb{R}^{B \times 4 \times 48 \times 48 \times 48}$.

  \item \textbf{Tangent (no grad):}
    $(\_, \tilde{v}) = f_\theta(z_t,\,r{=}t,\,t)$;
    $\tilde{v} \in \mathbb{R}^{B \times 4 \times 48 \times 48 \times 48}$
    is the v-head output at $r{=}t$ (instantaneous).
    Computed under \code{torch.no\_grad()}.

  \item \textbf{JVP (exact, dual):}
    $(u_\theta,\;du/dt,\;\hat{v}) =
    \texttt{jvp}(f_{\mathrm{dual}},\;(z_t,t,r),\;(\tilde{v},\mathbf{1},\mathbf{0}),\;
    \texttt{has\_aux})$.
    \begin{itemize}[nosep]
      \item $u_\theta \in \mathbb{R}^{B \times 4 \times 48 \times 48 \times 48}$
        (average velocity, derived from $\hat{x}$ via Eq.~\ref{eq:xpred}).
      \item $du/dt \in \mathbb{R}^{B \times 4 \times 48 \times 48 \times 48}$
        (JVP output).
      \item $\hat{v} \in \mathbb{R}^{B \times 4 \times 48 \times 48 \times 48}$
        (v-head output, auxiliary---no JVP through it).
    \end{itemize}

  \item \textbf{Compound velocity:}
    $V_\theta = u_\theta + (t-r).\mathrm{view}(B,1,1,1,1) \cdot \sg[du/dt]
    \in \mathbb{R}^{B \times 4 \times 48 \times 48 \times 48}$.

  \item \textbf{Per-sample raw losses:}
    $\ell_u^{(b)} = \|V_\theta^{(b)} - v_c^{(b)}\|_2^2 \in \mathbb{R}$,
    $\ell_v^{(b)} = \|\hat{v}^{(b)} - v_c^{(b)}\|_2^2 \in \mathbb{R}$;
    vectors $\boldsymbol{\ell}_u, \boldsymbol{\ell}_v \in \mathbb{R}^B$.

  \item \textbf{Adaptive weighting} (Eq.~\ref{eq:adaptive_simplified}):
    $\mathcal{L} = \frac{1}{B}\sum_{b=1}^{B}\!\left[
      \frac{\ell_u^{(b)}}{\sg[\ell_u^{(b)}]+1}
      + \frac{\ell_v^{(b)}}{\sg[\ell_v^{(b)}]+1}
    \right] \in \mathbb{R}$.

  \item \textbf{Backward + clip:}
    $\nabla_\theta \mathcal{L}$; clip
    $\|\nabla\|_2 \le 1.0$.

  \item \textbf{Optimiser step + EMA update:}
    $\theta \leftarrow \theta - \eta\,\nabla_\theta \mathcal{L}$;
    $\theta_\EMA \leftarrow \beta\,\theta_\EMA + (1-\beta)\,\theta$.
\end{enumerate}

\subsection{Time Sampling}
\label{sec:time_sampling}

Times are drawn from a logit-normal distribution.  Given a standard normal
sample $\xi \sim \mathcal{N}(0,1)$, the logit-normal transform is:
\begin{equation}
  t = \sigma(\mu + \sigma_t\,\xi)
    = \frac{1}{1 + e^{-(\mu + \sigma_t \xi)}}
    \;\in (0,1),
  \qquad
  t \leftarrow \max(t,\;t_{\min})
  \label{eq:time_sampling}
\end{equation}
where $\sigma(\cdot)$ is the sigmoid function, $\mu = -0.4$ (biases toward
lower $t$, \ie, near data), $\sigma_t = 1.0$, and $t_{\min} = 0.001$.
\begin{shapes}
$\xi$: $(B,)$; $t$: $(B,)$ after clamp.\\
Both $t$ and $r$ are sampled independently from the same distribution,
then swapped so $t \ge r$:\\
$t \leftarrow \max(t_\mathrm{raw}, r_\mathrm{raw})$;
$r \leftarrow \min(t_\mathrm{raw}, r_\mathrm{raw})$.\\
\textbf{Code:} \code{time\_sampler.py:62--73}.
\end{shapes}

\textbf{$\mathrm{data\_proportion} = 0.5$.}  For the first $\lfloor B/2
\rfloor$ samples in the batch, $r$ is set equal to $t$:
\begin{equation}
  r^{(b)} =
  \begin{cases}
    t^{(b)} & \text{if } b < \lfloor B/2 \rfloor
      \quad\text{(FM sample: }h = 0\text{, JVP vanishes)} \\
    r^{(b)}_{\mathrm{raw}} & \text{otherwise}
      \quad\text{(MF sample: }h > 0\text{, self-consistency enforced)}
  \end{cases}
  \label{eq:data_proportion}
\end{equation}

\subsection{Optimiser and Schedule}

Hyperparameters are listed in Table~\ref{tab:optimiser}.

\begin{table}[ht]
\centering
\caption{Optimiser hyperparameters.}
\label{tab:optimiser}
\begin{tabular}{lll}
\toprule
\textbf{Hyperparameter} & \textbf{Value} & \textbf{Justification} \\
\midrule
Optimiser       & AdamW         & Standard for diffusion/flow models \\
Learning rate   & $\eta = 10^{-4}$     & MeanFlow reference default \\
Weight decay    & 0.0           & No regularisation with adaptive loss \\
Betas           & $(\beta_1, \beta_2) = (0.9,\,0.95)$  & $\beta_2{=}0.95$: faster momentum adaptation \\
LR schedule     & Cosine decay  & $10^{-4} \to 5.6{\times}10^{-5}$ at early stop \\
Warmup          & 0 steps       & Warmup destabilised MF loss (Phase~4c) \\
Gradient clip   & $\|\nabla\|_2 \le 1.0$ & Prevents spike propagation \\
Mixed precision & bf16          & Standard for A100 \\
\bottomrule
\end{tabular}
\end{table}

\subsection{Compute Budget}

\begin{table}[ht]
\centering
\caption{Compute budget and actual training cost.}
\label{tab:compute}
\begin{tabular}{ll}
\toprule
\textbf{Property} & \textbf{Value} \\
\midrule
Hardware                   & 6$\times$ A100-SXM4-40GB (one DGX node) \\
Strategy                   & DDP (DistributedDataParallel) \\
Per-GPU batch              & $B_\mathrm{gpu} = 2$ (JVP activation: ${\sim}$20\,GB) \\
Gradient accumulation      & 11 steps \\
Effective batch            & $B_\mathrm{eff} = 2 \times 6 \times 11 = 132$ \\
Optimiser steps/epoch      & 41 \\
Planned optimiser steps    & 61,500 (1,500 epochs $\times$ 41) \\
Actual optimiser steps     & ${\sim}$28,980 (early-stopped at epoch~690) \\
Mean epoch time            & 268.1\,s (${\sim}$4.5\,min) \\
\textbf{Total wall time}   & \textbf{${\sim}$51.4\,hours (${\sim}$2.1\,days)} \\
\bottomrule
\end{tabular}
\end{table}

\subsection{EMA (Exponential Moving Average)}
\label{sec:ema}

An EMA model maintains shadow parameters $\theta_\EMA$ that are a running
weighted average of the online parameters $\theta$.  The update rule, applied
after each optimiser step, is:
\begin{equation}
  \theta_{\EMA}^{(k+1)}
    \coloneqq \beta\,\theta_{\EMA}^{(k)} + (1-\beta)\,\theta^{(k+1)}
  \label{eq:ema}
\end{equation}
where $\beta = 0.9999$ is the decay rate (half-life $\approx 6{,}931$ steps),
$k$ is the optimiser step index, $\theta^{(k+1)}$ are the parameters after the
$k$-th optimiser step, and $\theta_\EMA^{(0)} = \theta^{(0)}$ (initialised to
model weights).
\begin{shapes}
$\theta, \theta_\EMA$: all trainable parameters (${\sim}178$\,M scalars).\\
\textbf{Code:} \code{ema.py:50}: \code{shadow[name].mul\_(decay).add\_(param.data, alpha=1-decay)}.\\
Applied per-parameter: each shadow tensor has the same shape as its corresponding parameter.
\end{shapes}

\noindent
All FID evaluations during training use EMA weights.  The best model checkpoint
is selected by FID computed from EMA-weight samples.

\subsection{Divergence Monitoring}

An EMA-smoothed raw loss monitor (decay~0.99, half-life ${\sim}$69 steps)
tracks training stability.  Warnings fire at $3\times$ and $5\times$ the
minimum EMA loss, with a grace period of 1000~steps.  Early stopping is
FID-based via \code{EvaluationCallback} with patience\,=\,10.


%═══════════════════════════════════════════════════════════════════════════════
\section{Sampling and Generation}
\label{sec:sampling}

\subsection{One-Step Sampling (1-NFE)}

With x-prediction, 1-NFE sampling is a single forward pass:
\begin{equation}
  \varepsilon \sim \mathcal{N}(\mathbf{0},\,I_d)
    \;\in \mathbb{R}^{B \times C \times D \times H \times W},
  \qquad
  z_0 = f_\theta(\varepsilon,\;r{=}\mathbf{0}_B,\;t{=}\mathbf{1}_B)
    \;\in \mathbb{R}^{B \times C \times D \times H \times W}
  \label{eq:1nfe_sampling}
\end{equation}
\begin{shapes}
$\varepsilon, z_0$: $(B, 4, 48, 48, 48)$.
\quad $r = \mathbf{0}_B$, $t = \mathbf{1}_B$: each $(B,)$.\\
The v-head output is discarded: \code{if isinstance(output, tuple): output = output[0]}.\\
\textbf{Code:} \code{one\_step.py:32--46}.
\end{shapes}

\subsection{Multi-Step Euler Sampling}

For comparison, multi-step Euler sampling divides $[1,0]$ into $K$ uniform
steps with timesteps $t_i = 1 - i/K$ for $i = 0, \ldots, K$:
\begin{equation}
  \Delta t_i \coloneqq t_i - t_{i+1} = \frac{1}{K},
  \qquad
  t_i, t_{i+1} \in \{1,\;1{-}\tfrac{1}{K},\;\ldots,\;0\}
  \label{eq:euler_times}
\end{equation}
At each step, the model predicts $\hat{x}$, which is converted to velocity
and used for an Euler update:
\begin{equation}
  \hat{x}_i = f_\theta(z^{(i)},\;r{=}t_{i+1},\;t{=}t_i)
    \;\in \mathbb{R}^{C \times D \times H \times W},
  \qquad
  u_i = \frac{z^{(i)}-\hat{x}_i}{\max(t_i,\,0.05)},
  \qquad
  z^{(i+1)} = z^{(i)} - \Delta t_i \cdot u_i
  \label{eq:euler_step}
\end{equation}
\begin{shapes}
$z^{(i)}, \hat{x}_i, u_i$: each $(B, 4, 48, 48, 48)$.
\quad $z^{(0)} = \varepsilon$; $z^{(K)} = z_0$ (final sample).\\
\textbf{Code:} \code{multi\_step.py}; NFE levels tested: $K \in \{1,2,5,10\}$.
\end{shapes}

\subsection{Full Generation Pipeline (Phase~5)}

\begin{enumerate}[nosep]
  \item Pre-generate shared noise $\varepsilon \in \mathbb{R}^{N \times 4 \times 48 \times 48 \times 48}$ (for NFE-consistency analysis).
  \item For each NFE level: generate latents, store in HDF5 (latents, seeds,
        timing).
  \item Denormalise: $z_0 = \hat{z}\cdot\sigma_{\mathrm{latent}}
                           + \mu_{\mathrm{latent}}$,
    where $\mu_\mathrm{latent}, \sigma_\mathrm{latent} \in \mathbb{R}^{1 \times 4 \times 1 \times 1 \times 1}$.
  \item Decode: $\hat{x}_\mathrm{MRI} = \mathrm{Dec}_\phi(z_0\,/\,s)
    \in \mathbb{R}^{B \times 1 \times 192 \times 192 \times 192}$.
  \item Clamp to $[0,1]$; store decoded volumes in HDF5.
\end{enumerate}

\textbf{Shared noise protocol.}  The same noise tensor $\varepsilon$ is used
across all NFE levels for a given sample index, enabling direct comparison of
how additional steps refine the output.


%═══════════════════════════════════════════════════════════════════════════════
\section{Experimental Results}
\label{sec:results}

\subsection{Toy Validation (Phase~2): MeanFlow on Flat Torus}

Before scaling to brain MRI, MeanFlow was validated on a flat torus in
$\mathbb{R}^4$ (Table~\ref{tab:torus}).

\begin{table}[ht]
\centering
\caption{Phase~2 toy validation results.}
\label{tab:torus}
\begin{tabular}{llr}
\toprule
\textbf{Metric} & \textbf{Result} & \textbf{Target} \\
\midrule
1-NFE torus distance            & 0.0892  & $<$0.1 \\
5-step torus distance           & 0.0271  & $<$0.05 \\
Angular KS $p$-value ($\theta_1$) & 0.374 & $>$0.01 \\
Angular KS $p$-value ($\theta_2$) & 0.575 & $>$0.01 \\
x-pred vs u-pred loss ratio     & $<$1.2  & $<$1.5 \\
\bottomrule
\end{tabular}
\end{table}

This confirmed: (i)~the MeanFlow implementation correctly learns 1-step
generation, (ii)~the angular distribution is statistically indistinguishable
from truth (KS test), and (iii)~both x-prediction and u-prediction converge.

\subsection{Main Training (Phase~4): Best Model}

The best model uses x-prediction + exact JVP + $(t,h)$ conditioning + v-head
(1~ResBlock) + $L_2$ loss (Table~\ref{tab:best_model}).

\begin{table}[ht]
\centering
\caption{Best model configuration.}
\label{tab:best_model}
\begin{tabular}{ll}
\toprule
\textbf{Property} & \textbf{Value} \\
\midrule
Prediction type   & x-prediction \\
JVP strategy      & Exact (\code{torch.func.jvp}) \\
Conditioning      & $t\_h$ (condition on $t$ and $h=t{-}r$) \\
v-head            & 1 ResBlock (${\sim}$228\,K params, $<$0.13\%) \\
Loss norm         & $L_2$ ($p{=}2$) with independent adaptive weighting \\
Best FID epoch    & \textbf{388} (step 16,338) \\
Early-stopped at  & Epoch 690 (patience=10, 46\% of planned 1,500) \\
\bottomrule
\end{tabular}
\end{table}

\subsection{FID Progression}

FID is computed as a 2.5D metric using RadImageNet ResNet-50 features on
central axial, coronal, and sagittal slices (Table~\ref{tab:fid}).

\begin{table}[ht]
\centering
\caption{2.5D FID progression during training (selected checkpoints).}
\label{tab:fid}
\begin{tabular}{rrrrrl}
\toprule
\textbf{Epoch} & \textbf{Step} & $\FID_{\mathrm{avg}}$ & $\FID_{xy}$ & $\FID_{yz}$ & \textbf{Notes} \\
\midrule
9    & 420    & 46.76 & 45.12 & 46.68 & First evaluation \\
28   & 1,218  & 28.40 & 17.25 & 43.63 & Rapid initial drop \\
88   & 3,738  & 14.61 & 14.08 & 16.25 & Approaching plateau \\
178  & 7,518  & 14.06 & 13.79 & 15.46 & \\
268  & 11,298 & 12.94 & 12.77 & 14.06 & \\
\textbf{388} & \textbf{16,338} & \textbf{11.67} & \textbf{11.85} & \textbf{12.09} & \textbf{Best FID} \\
628  & 26,418 & 11.92 & 12.03 & 12.33 & \\
688  & 28,938 & 11.88 & 12.14 & 12.08 & Final evaluation \\
\bottomrule
\end{tabular}
\end{table}

FID drops rapidly from 46.76 to ${\sim}$14 in the first 88~epochs, then
gradually improves to \textbf{11.67} at epoch~388, after which it plateaus.
Training continued for 302~more epochs without improvement before early
stopping triggered.

\subsection{Training Dynamics}

Table~\ref{tab:dynamics} shows key training metrics at milestone epochs.

\begin{table}[ht]
\centering
\caption{Training dynamics at milestone epochs.}
\label{tab:dynamics}
\begin{tabular}{rrccccc}
\toprule
\textbf{Epoch} & \textbf{Raw Loss} & $\cos(V,v_c)$ & $\cos(\hat{v},v_c)$ & \textbf{Rel.\,Err.} & $\|\nabla\|$ & \textbf{LR} \\
\midrule
0   & 2,788,299 & 0.083 & 0.170 & 1.353 & 1.384 & $1.0\!\times\!10^{-4}$ \\
50  & 5,316,270 & 0.240 & 0.388 & 1.759 & 0.953 & ${\sim}9.9\!\times\!10^{-5}$ \\
100 & 6,275,983 & 0.274 & 0.412 & 1.778 & 0.910 & ${\sim}9.7\!\times\!10^{-5}$ \\
388 & 4,724,751 & 0.295 & 0.429 & 1.645 & 1.180 & $8.5\!\times\!10^{-5}$ \\
689 & 4,327,238 & 0.307 & 0.436 & 1.516 & 0.610 & $5.6\!\times\!10^{-5}$ \\
\bottomrule
\end{tabular}
\end{table}

\textbf{Key observations:}
\begin{enumerate}[nosep]
  \item \textbf{Raw loss is not a good proxy for generation quality.}  The raw
    loss is dominated by the MF term ($\ell_{\mathrm{MF}} \approx
    6.5\!\times\!10^6$ vs $\ell_{\mathrm{FM}} \approx
    7\!\times\!10^5$) and does not monotonically decrease.  FID (computed from
    EMA-weight samples) is the correct early-stopping metric.
  \item \textbf{Cosine alignment improves steadily.}  $\cos(V,v_c)$ increases
    from 0.083 to 0.307; the v-head alignment converges faster (0.170 $\to$
    0.436), confirming it provides good tangents from early training.
  \item \textbf{Gradient norm decreases dramatically.}  From 1.384 to 0.610;
    clipping fraction drops ${\sim}90\%\!\to\!{\sim}5\%$, indicating stable
    dynamics by epoch~${\sim}300$.
\end{enumerate}

\subsection{Generated Sample Statistics}

The model's output $\hat{x}$ statistics evolve during training
(Table~\ref{tab:xhat}).

\begin{table}[ht]
\centering
\caption{Statistics of generated $\hat{x} \in \mathbb{R}^{4 \times 48 \times 48 \times 48}$ at different training stages.}
\label{tab:xhat}
\begin{tabular}{rcccc}
\toprule
\textbf{Epoch} & $\mathbb{E}[\hat{x}]$ & $\mathrm{std}[\hat{x}]$ & $\min[\hat{x}]$ & $\max[\hat{x}]$ \\
\midrule
0   & $+0.004$ & $0.243$ & $-2.06$ & $3.41$ \\
100 & $-0.001$ & $0.737$ & $-5.88$ & $8.56$ \\
388 & $-0.001$ & $0.753$ & $-5.84$ & $8.69$ \\
689 & $-0.001$ & $0.781$ & $-6.22$ & $9.38$ \\
\bottomrule
\end{tabular}
\end{table}

The standard deviation stabilises around 0.75--0.78, approximately 75\% of the
true latent std (${\sim}$0.97--1.02).  This \textbf{variance under-estimation}
is the primary quality bottleneck: generated samples have slightly less contrast
than real ones.

\subsection{NFE Quality Analysis}

Visual inspection of decoded samples reveals a clear quality hierarchy
(Table~\ref{tab:nfe_quality}).

\begin{table}[ht]
\centering
\caption{Visual quality at different NFE levels.}
\label{tab:nfe_quality}
\begin{tabular}{lll}
\toprule
\textbf{NFE} & \textbf{Quality} & \textbf{Key Features} \\
\midrule
1  & Recognisable, blurry    & Global structure correct, low contrast \\
2  & Substantially sharper   & Cortical sulci visible \\
5  & Sharp, plausible        & Clear grey-white contrast \\
10 & Marginally crisper      & Diminishing returns beyond 5 \\
\bottomrule
\end{tabular}
\end{table}

\textbf{NFE-consistency metrics} (shared noise, converged model):
cosine similarity NFE=1 vs NFE=2: ${\sim}$0.96; vs NFE=5: ${\sim}$0.91;
vs NFE=10: ${\sim}$0.90.

The 1-NFE output captures global anatomy correctly but lacks high-frequency
detail---a known MeanFlow limitation.  NFE=2 provides substantial improvement
at only $2\times$ cost.

\subsection{Spectral Analysis}

The radially-averaged power spectrum evolves from flat (noise-like at epoch
${\sim}100$) to the characteristic $1/f$ brain spectrum by epoch~600.  All
4~channels show proper low-frequency dominance with high-frequency rolloff.

\subsection{SWD (Sliced Wasserstein Distance)}

SWD improves rapidly in early training (best = 1.779 at epoch~9) but then
\textbf{diverges from FID}---SWD increases monotonically after epoch ${\sim}50$
(final: 2.555) while FID continues to decrease.  SWD is \textbf{not a reliable
proxy} for perceptual quality in this setting.

\subsection{Key Ablation Insights}

\textbf{x-pred vs u-pred.}  x-prediction with exact JVP trained stably through
690~epochs (best FID at~388).  This supports the pMF manifold hypothesis:
$d_{\mathrm{input}}/d_{\mathrm{bottleneck}} \approx 4$, firmly in the
x-prediction regime.

\textbf{x-pred + FD-JVP instability.}  Exponential growth of JVP norms
($10{,}000\times$ within 50~epochs).  Root cause (Eq.~\ref{eq:fd_instability}):
\begin{equation}
  u = \frac{z_t - \hat{x}}{t} \implies
  \frac{du}{dt}\bigg|_{\mathrm{FD}}
    \approx \frac{u(t{+}\delta)-u(t)}{\delta}
    \sim \frac{O(1/t)}{\delta}
    = O\!\left(\frac{1}{t^2 \delta}\right)
\end{equation}

\textbf{Conditioning mode.}  $(t,h)$ selected per MeanFlow Table~1c: FID~61.06
for $(t,h)$ vs~63.13 for $h$-only.

\textbf{Training stability fixes (Phase~4c--4e):}
\begin{itemize}[nosep]
  \item $\beta_2$: 0.999 $\to$ 0.95 (prevents stale gradients).
  \item Warmup: 5000 $\to$ 0 (destabilises MF loss).
  \item $\varepsilon_{\mathrm{norm}}$: 0.01 $\to$ 1.0 (prevents runaway
        amplification).
  \item $p_{\mathrm{norm}}$: 0.5 $\to$ 1.0 (sub-linear normalisation
        $\to$ gradient explosion).
  \item $\mathrm{data\_proportion}$: stabilised at 0.5 (50/50 FM/MF split).
\end{itemize}


%═══════════════════════════════════════════════════════════════════════════════
\section{Evaluation Protocol}
\label{sec:evaluation}

\subsection{Image Quality Metrics}

\begin{table}[ht]
\centering
\caption{Evaluation metrics suite.}
\label{tab:metrics}
\begin{tabular}{lll}
\toprule
\textbf{Metric} & \textbf{Description} & \textbf{Features} \\
\midrule
2.5D FID  & Slice-wise FID (ax/cor/sag)   & RadImageNet ResNet-50 \\
3D FID    & Volumetric FID               & Med3D ResNet-50 \\
MMD       & Max Mean Discrepancy         & Gaussian kernel \\
Coverage  & Real--synthetic NN fraction  & Feature-space NN \\
Density   & Synthetic near real manifold & Complementary to coverage \\
MS-SSIM   & 3-level Multi-Scale SSIM (3D) & MONAI \code{SSIMMetric} \\
PSNR      & Peak Signal-to-Noise Ratio   & Actual data range \\
\bottomrule
\end{tabular}
\end{table}

\subsection{Spectral Analysis}

High-frequency energy ratio via 3D FFT:
\begin{equation}
  \rho \coloneqq \frac{\displaystyle\sum_{|\mathbf{k}|>k_0} |\mathcal{F}(x)(\mathbf{k})|^2}
              {\displaystyle\sum_{\mathbf{k}} |\mathcal{F}(x)(\mathbf{k})|^2}
    \;\in [0,1]
  \label{eq:spectral}
\end{equation}
where $\mathcal{F}(x) \in \mathbb{C}^{D' \times H' \times W'}$ ($D'{=}H'{=}W'{=}192$)
is the 3D discrete Fourier transform, $\mathbf{k} \in \mathbb{Z}^3$ is the
frequency index, $|\mathbf{k}|$ is the radial frequency, and $k_0$ is the
cutoff.  Reported for real, VAE-reconstructed, and generated volumes.

\subsection{Morphological Assessment}

SynthSeg~\cite{synthseg} segmentation on both real and synthetic volumes:
regional volume correlation, distributional KL divergence, and SynthSeg Dice
overlap (paired by nearest-neighbour in feature space).

\subsection{NFE Consistency Analysis}

Generated samples at NFE $\in \{1,2,5,10\}$ using shared noise, compared via
per-sample $L_2$ distance and FID at each NFE level.


%═══════════════════════════════════════════════════════════════════════════════
\section{Implementation Details}
\label{sec:implementation}

\subsection{Software Stack}

Python 3.11.14, PyTorch 2.10+cu128, MONAI 1.5.2, PyTorch Lightning 2.6.1,
OmegaConf/Hydra, HDF5 (h5py), Python logging with Rich handler + TensorBoard.

\subsection{Gated Phase System}

The project is implemented in 9~gated phases (0--8).  Phase $N{+}1$ cannot
begin until Phase $N$'s critical tests all pass (Table~\ref{tab:phases}).

\begin{table}[ht]
\centering
\caption{Phase status.}
\label{tab:phases}
\begin{tabular}{clll}
\toprule
\textbf{Phase} & \textbf{Title} & \textbf{Status} & \textbf{Tests} \\
\midrule
0 & VAE Validation         & Complete      & 7/7 \\
1 & Latent Pre-computation & Complete      & 7/7 \\
2 & Toy Experiment (Torus) & Complete      & 6/6 + 2 info \\
3 & MeanFlow Loss + 3D UNet & Complete    & 28/28 \\
4 & Training on Brain MRI  & Complete      & 12+ \\
5 & Generation + Evaluation & Code complete & 11/11 local \\
6 & Ablation Runs          & Planned       & --- \\
7 & LoRA Fine-Tuning (FCD) & Planned       & --- \\
8 & Paper Figures          & Planned       & --- \\
\bottomrule
\end{tabular}
\end{table}

\subsection{Compute Infrastructure}

\begin{itemize}[nosep]
  \item \textbf{Local:} RTX 4060 8\,GB --- development, testing, analysis.
  \item \textbf{Picasso:} 4 nodes $\times$ 8$\times$ A100-40GB --- training,
        encoding, evaluation.  SLURM with \code{--constraint=dgx}.
\end{itemize}


%═══════════════════════════════════════════════════════════════════════════════
\section{Novel Contributions}
\label{sec:contributions}

\textbf{Contribution 1: First MeanFlow for 3D Medical Image Synthesis.}
We achieve 1-NFE generation of $192^3$ brain MRI volumes in the latent space
of a frozen MAISI VAE---a \textbf{50--1000$\times$ reduction in sampling cost}
compared to existing methods (DDPM: 1000; flow matching: 10--50; rectified
flow: 5--50 steps).  Best 2.5D FID: \textbf{11.67}, trained in ${\sim}$51
hours on 6$\times$ A100 GPUs.

\textbf{Contribution 2: Per-Channel $L_p$ Loss in Latent MeanFlow.}
We extend the SLIM-Diff per-channel $L_p$ loss~\cite{slimdiff} from
pixel-space DDPM to latent-space MeanFlow, providing a framework for
ablation over $p \in \{1.0,\ldots,3.0\}$.

\textbf{Contribution 3: x-Prediction vs u-Prediction in Latent 3D UNets.}
First investigation of this dichotomy (established by pMF for 2D ViTs) in
latent space with 3D UNets.  We document the critical incompatibility:
\textbf{x-prediction + FD-JVP is numerically unstable} due to the $1/t$
singularity.

\textbf{Contribution 4: iMF Dual-Head in Medical Context.}
First application of the iMF dual-head architecture outside natural images.
The v-head adds ${\sim}$228\,K parameters ($<$0.13\%) and is disabled at
inference.

\textbf{Contribution 5 (Planned): LoRA Fine-Tuning for Joint Image-Mask
Synthesis.}  Phase~7 will demonstrate LoRA fine-tuning for joint FLAIR MRI +
FCD mask synthesis in a data-scarce setting (${\sim}$50--100 cases).


%═══════════════════════════════════════════════════════════════════════════════
\section{Current Status and Next Steps}
\label{sec:status}

\subsection{Current State (February 2026)}

\begin{itemize}[nosep]
  \item \textbf{Phases 0--4:} Complete. Best model trained (x-pred + exact JVP,
    best FID 11.67 at epoch~388, early-stopped at epoch~690).
  \item \textbf{Phase 5:} Code complete, 11/11 local tests pass.  Awaiting
    Picasso execution for latent generation at NFE$\in\{1,2,5,10\}$, VAE
    decoding, and full quantitative evaluation.
  \item \textbf{Phases 6--8:} Planned.
\end{itemize}

\subsection{Key Findings}

\begin{enumerate}[nosep]
  \item \textbf{Early stopping was effective:} Best FID at epoch~388,
    used only 46\% of planned 1,500~epochs. Total training: 51.4~hours.
  \item \textbf{1-NFE quality gap is significant:} NFE=1 is blurry;
    NFE$\geq$2 is substantially sharper.
  \item \textbf{Variance under-estimation:} Generated std ${\sim}$0.78 vs
    true ${\sim}$0.97--1.02.  The 20--25\% gap is the primary quality
    bottleneck.
  \item \textbf{v-head converges fast:} $\cos(\hat{v},v_c)$ reaches~0.39 by
    epoch~50 (85\% of final~0.44).
  \item \textbf{SWD anti-correlates with FID} beyond early training.
  \item \textbf{Gradient clipping fraction:} 90\% $\to$ 5\%, confirming stable
    dynamics by epoch~${\sim}300$.
\end{enumerate}

\subsection{Open Questions}

\begin{itemize}[nosep]
  \item Does the $L_p$ exponent effect from pixel space transfer to latent
    space?
  \item How much does VAE smoothing degrade MeanFlow volumes vs multi-step
    methods?
  \item Can LoRA fine-tuning on ${\sim}$50--100 FCD cases produce anatomically
    plausible lesion synthesis?
  \item Can the variance under-estimation (${\sim}$75\% of true std) be
    mitigated by post-hoc rescaling or loss modifications?
\end{itemize}


%═══════════════════════════════════════════════════════════════════════════════
\begingroup
\renewcommand{\section}[2]{}  % suppress natbib's \section*{References}
\begin{thebibliography}{15}

\bibitem{meanflow}
Z.~Geng, M.~Deng, X.~Bai, J.\,Z.~Kolter, and K.~He,
``Mean Flows for One-step Generative Modeling,''
\emph{NeurIPS}, 2025 (Oral). arXiv:2505.13447.

\bibitem{imf}
Z.~Geng, Y.~Lu, Z.~Wu, E.~Shechtman, J.\,Z.~Kolter, and K.~He,
``Improved Mean Flows: On the Challenges of Fastforward Generative Models,''
arXiv:2512.02012, 2025.

\bibitem{pmf}
Y.~Lu, S.~Lu, Q.~Sun, H.~Zhao \etal,
``One-step Latent-free Image Generation with Pixel Mean Flows,''
arXiv:2601.22158, 2026.

\bibitem{maisi}
P.~Guo, C.~Zhao, D.~Yang \etal,
``MAISI: Medical AI for Synthetic Imaging,''
\emph{WACV}, 2025. arXiv:2409.11169.

\bibitem{maisi_v2}
C.~Zhao, P.~Guo, D.~Yang \etal,
``MAISI-v2: Accelerated 3D High-Resolution Medical Image Synthesis with
Rectified Flow,''
arXiv:2508.05772, 2025.

\bibitem{slimdiff}
M.~Pascual-Gonz\'{a}lez \etal,
``SLIM-Diff: Shared Latent Image-Mask Diffusion with Lp Loss for Data-Scarce
Epilepsy FLAIR MRI,''
arXiv:2602.03372, 2026.

\bibitem{motfm}
M.~Yazdani \etal,
``Flow Matching for Medical Image Synthesis,''
\emph{MICCAI}, 2025. arXiv:2503.00266.

\bibitem{medddpm}
Z.~Dorjsembe \etal,
``Semantic 3D Brain MRI Synthesis with Channel-Wise Conditioning,''
\emph{IEEE JBHI}, 2024.

\bibitem{lipman2023}
Y.~Lipman, R.\,T.\,Q.~Chen, H.~Ben-Hamu, M.~Nickel, and M.~Le,
``Flow Matching for Generative Modeling,''
\emph{ICLR}, 2023.

\bibitem{liu2023}
X.~Liu, C.~Gong, and Q.~Liu,
``Flow Straight and Fast: Learning to Generate and Transfer Data with Rectified
Flow,''
\emph{ICLR}, 2023.

\bibitem{ldm}
R.~Rombach, A.~Blattmann, D.~Lorenz, P.~Esser, and B.~Ommer,
``High-Resolution Image Synthesis with Latent Diffusion Models,''
\emph{CVPR}, 2022.

\bibitem{lora}
E.~Hu, Y.~Shen, P.~Wallis \etal,
``LoRA: Low-Rank Adaptation of Large Language Models,''
\emph{ICLR}, 2022.

\bibitem{synthseg}
B.~Billot \etal,
``SynthSeg: Segmentation of Brain MRI Scans of Any Contrast and Resolution,''
\emph{Medical Image Analysis}, 2023.

\bibitem{ddpm}
J.~Ho, A.~Jain, and P.~Abbeel,
``Denoising Diffusion Probabilistic Models,''
\emph{NeurIPS}, 2020.

\bibitem{brlp}
R.~Puglisi \etal,
``Brain Latent Progression (BrLP),''
\emph{Medical Image Analysis}, 2025.

\end{thebibliography}
\endgroup

\end{document}
